\chapter{Fizyczne podstawy}
% =======================================================
\section{Dynamika bryły sztywnej w 6 stopniach swobody (6-DOF)}
% Teoria: Masa, inercja, środek ciężkości. Kąty Eulera: Roll (Przechył), Pitch (Pochylenie), Yaw (Odchylenie). Jak wektor ciągu silników wpływa na ruch w konkretnym kierunku.

\section{Zakłócenia zewnętrzne i estymacja stanu}
% Teoria: Szum czujników (GPS, barometr, IMU). Dlaczego pomiary nigdy nie są idealne?
\subsection{Zjawisko znoszenia (Wind drift) i Rozszerzony Filtr Kalmana (EKF)}
% Teoria: Fuzja sensorów. Jak dron domyśla się obecności wiatru analizując różnicę między modelem matematycznym a faktycznym przesunięciem GPS.
\subsection{Przykład praktyczny: Nasłuchiwanie telemetrii w warunkach silnego wiatru}
% Kod Python: Skrypt 04_physics_wind.py. Odczyt danych z MAVLinka i analiza kompensacji wektora wiatru.

